\section{Core Components}
\subsection{Bootstrapping, MMIO, and Interrupt Setup}

The kernel starts without support from a host OS or a C runtime. The first executed project code is AArch64 assembly. It sets up the initial stack, clears the \texttt{.bss} section, and then branches into the C kernel entry point. Clearing \texttt{.bss} is required because global kernel objects such task tables, mutexes, device state, and buffers must start in a known zero-initialized state.

After entering C code, the kernel initializes the hardware-facing subsystems before starting the scheduler. UART is initialized early to provide debug output. The interrupt controller is configured for the interrupts used by the kernel: the mini UART receive interrupt, the generic timer interrupt, the GPIO interrupt used by the Sense HAT joystick, and the I2C interrupt. The IRQ dispatcher acknowledges the interrupt at the GIC CPU interface, dispatches to the responsible subsystem, and then signals end-of-interrupt back to the controller.

\subsection{Task System and Cooperative Scheduler}

The task system is based on a statically allocated task table. Each \textbf{task control block (TCB)} contains the task ID, state, flags, saved stack pointer, entry function, wakeup tick, name and private stack. Dynamic allocation is deliberately avoided for TCBs and stacks, so the maximum number of tasks and the memory cost of each task are known at compile time. Task creation prepares the TCB and constructs an artificaial initial stack frame. This frame matches the layout expected by the AArch64 context switch routine. The saved return address points to \texttt{task\_bootstrap()}. When the scheduler selects the task for the first time, the context switch restores this prepared context and returns into the bootstrap function. The same low-level context switch mechanism is therefore used both for starting a new task and for resuming an existing task.

The \textbf{scheduler} is cooperative. A task runs until it explicitly yields, sleeps, blocks, or exits. The scheduler scans the task table in round-robin order and selects the next ready non-idle task. If no ordinary task is ready, the idle-task is selected. The idle task executes a wait-for-event instruction and yields again, so the CU is not kept in a busy loop while no task can make progress. The task state encode why a task is or is not runnable:

\begin{itemize}
	\item A \texttt{READY} task may be selected by the scheduler.
	\item A \texttt{RUNNING} task currently owns the CPU.
	\item A \texttt{SLEEPING} task is reactivated by the timer interrupt once a tick is reached.
	\item A \texttt{BLOCKED} task waits for an event, such as a mutex becoming available or a device interrupt.
	\item A \texttt{DYING} task is no longer scheduled and is cleaned up by the scheduler at a safe point.
\end{itemize}

Not all tasks are driven the same way. Some tasks are purely cooperative: they perform, work, sleep for a number of ticks, and then continue. Examples include periodic demo tasks, dashboards, and visualization. Their timing is based on explicit calls to \texttt{task\_sleep()}.

Other tasks are interrupt-driven. They normally block and are only made ready when an interrupt indicates that work is available. The joystick task is the clearest example. The GPIO interrupt handler detects the Sense HAT joystick interrupt line, clears the GPIO event, and wakes the joystick task. The task then runs outside interrupt context, reads the joystick state over I2C, converts the hardware state into logical menu events, and dispatches them to the menu code. Sleeping tasks are also connected to interrupt handling. The generic timer interrupt increments the software tick counter and checks all sleeping tasks. If a task's wakeup tick has reached, the timer code changes its state back to \texttt{READY}. The timer interrupt does not preemptively switch tasks but only updates the scheduler-visible state.

\subsection{Synchronization and Shared Hardware Access}

Even with cooperative scheduling, \textbf{synchronization} is required because multiple tasks can acces shared kernel objects and because interrupts may update state asynchronously. The kernel therefore uses short critical sections where interrupts are disabled while sensitive data structures are modified. This is used for operations such as changing task state, waking tasks, updating scheduler data, or manipulating synchronization primitves.

For longer shared-resource protection, the kernel provides \textbf{mutexes}. A mutex records whether it is locked and which task owns it. If another task tries to acquire the mutex while it is already locked, the task is put into a wait queue, marked as blocked, and the scheduler selects another runnable task. When the mutex is released, one waiting task is woken up and can continue later.

The most important protected shared resource is the I2C bus. Several independent components compete for the same bus:

\begin{itemize}
	\item The environmental task reads pressure, humidity, and temperature sensors.
	\item The joystick task reads Sense HAT joystick state after GPIO interrupt.
	\item LED-related tasks update the Sense HAT LED matrix.
	\item Dashboard and shell commands may request recent sensor data or trigger display updates.
\end{itemize}

These operations must not interleave because an I2C transaction consists of multiple ordered register accesses. The global I2C bus mutex serializes these transactions. Every driver-level access to a Sense HAT device must acquire the bus lock before starting communication and release it afterwards. This meas the synchronization boundary is not the individual sensor or LED matrix, but the bus itself. Protecting the bus prevents one task from changing the I2C target address, FIFO state, or transfare control while another transaction is in progress.

A similar ownership principle is used for display output. HDMI output is divides into panes, and tasks can acquire exclusive ownership of a pane before rendering application-specific output. The LED matrix is also treated as a shared hardware resource. Instead of allowing arbitrary concurrent writes, LED output is routed through controlled rendering code so that hardware access remains centralized.

\subsection{HDMI Rendering}

HDMI output is one of the more complex parts of the project. The kernel requests a framebuffer from the Raspberry Pi firmware through the mailbox property interface. After initialization, rendering is performed directly by writing pixels into this framebuffer. Since the kernel does not use a graphical library or terminal emulator, text output, scrolling, clearing, and application views all have to be implemented manually.

The final design does not redraw the screen line by line on every update. Instead, HDMI output is modeled as a \textbf{character-cell display with backing state}. Each pane stores the character and color information for its cells. When code writes to a pane, it updates the cell state and marks only the affected cells as dirty. A later presentation step flushes dirty cells by converting their glyphs into framebuffer pixels. The dirty-cell model is important for responsiveness. A full redraw of the HDMI framebuffer is expensive and would let display output dominate CPU time and disturb the cooperative scheduler. Dirty-cell flushing limits the amount of work per update and makes it possible to budget rendering work across scheduler iterations.

\subsection{Sense HAT LED Matrix Rendering}

The Sense HAT LED matrix is also treated as shared output device. The matrix has 64 logical pixels, but the device expects color data to be transferred through the Sense HAT I2C interface. The kernel therefore builds LED frames in memory first and then sends the complete frame to the matrix controller. The LED matrix competes with the environmental sensors and joystick handling for the same I2C bus. Therefore, presenting a LED frame also has to respect the global I2C lock. This is especially relevant because frame updates are larger than simple sensor register reads: sending a full LED frame means transferring many color bytes, so it must be treated as an exclusive bus operation.

\subsection{Diagnostics}

Because the system runs without Linux, there are no external tools such as \texttt{top}, \texttt{dmesg}, or a process monitor. The kernel therefore implements its own diagnostics. The shell can print scheduler state, task states, heap usage, stack high-water marks, logs, and trace entries. Since log entries are heap-allocated, heap diagnostics also show whether runtime logging remains bounded. The diagnostics dashboard renders selected runtime data live to the HDMI main pane.
